位运算和状态压缩深度题卡 按题型拆成章内大节，但每个大节都先从 LeetCode 案例切入，再进入分流、案例矩阵、案例推导和实验复盘。这样读者看到的是从具体题推到规律，而不是先读抽象方法再猜它能用在哪。

\topic{位运算与状态压缩}
这一节先看 LeetCode 29 两数相除、LeetCode 50 Pow(x, n)、LeetCode 136 只出现一次的数字 这几道 LeetCode 题，再整理同一题型的分流和推导。阅读顺序不是先背原理，而是先看案例的暴力瓶颈，再把重复出现的观察抽象成方法。

\subsection{原理和适用条件}

以 LeetCode 29 两数相除、LeetCode 50 Pow(x, n)、LeetCode 136 只出现一次的数字 为主线：LeetCode 29 两数相除：模型是“除法可以看成不断减去除数的倍增值。”，优化抓手是“把除数按二进制倍增，优先减去不超过被除数的最大倍数并累加商。”；LeetCode 50 Pow(x, n)：模型是“指数 n 可以按二进制拆分，重复乘 x 的线性过程存在大量倍增结构。”，优化抓手是“快速幂每轮根据最低位决定是否乘入答案，同时底数平方、指数右移。”；LeetCode 136 只出现一次的数字：模型是“成对元素可以用异或抵消，剩下单个元素。”，优化抓手是“异或满足交换结合，扫描顺序不影响答案。”。这些案例共同推出的规律是：先用计数或哈希保证正确，再寻找位运算的代数抵消、按位统计或函数图结构。位题的难点是证明被抵消的信息确实不影响答案。 方法名只是复盘后的标签，真正要掌握的是从题面约束、暴力失败、状态设计到边界反例的推导链。

\begin{principlebox}{图示：从 LeetCode 案例到 位运算与状态压缩推导链}
\begin{tabular}{@{}p{0.18\linewidth}p{0.74\linewidth}@{}}
暴力基线 & 枚举候选、完整模拟或逐个检查，先保证覆盖全部可能答案。\\
瓶颈定位 & 慢点在于状态复制和集合比较。位操作让包含、加入、删除、枚举子集变成常数或接近常数的机器指令。\\
结构选择 & 异或解决成对抵消，按位计数处理出现三次，掩码表示访问集合或可用集合，低位提取用于枚举候选。\\
不变量 & 掩码的每一位必须有固定含义。状态去重时，位置相同但掩码不同不能合并，因为未来能力不同。\\
代码落地 & 使用无符号整数能减少移位歧义。位数超过 31 时用 \code{long long} 或 \code{std::bitset}。表达式加括号，避免位运算优先级误读。\\
\end{tabular}
\end{principlebox}

\subsection{LeetCode 案例分流}

\begin{principlebox}{从 LeetCode 案例分流：位运算与状态压缩}
\textbf{异或抵消。}
代表案例：LeetCode 136 只出现一次的数字、LeetCode 260 只出现一次的数字 III。先看这些题的题面条件：除一个或两个数外，其他数按偶数次出现。由此推出的处理方式是：利用异或交换结合和 x 异或 x 为零。风险点：出现次数不是两次时不能直接异或。

\textbf{按位计数。}
代表案例：LeetCode 137 只出现一次的数字 II、LeetCode 剑指 Offer 56-I 数组中数字出现次数。先看这些题的题面条件：其他数字出现固定 k 次，目标数字出现不同次数。由此推出的处理方式是：逐位统计一的数量，对 k 取模还原答案。风险点：负数和符号位要使用足够宽的整数。

\textbf{掩码集合。}
代表案例：LeetCode 78 子集、LeetCode 90 子集 II、LeetCode 864 获取所有钥匙的最短路径。先看这些题的题面条件：状态是若干布尔选择或访问集合。由此推出的处理方式是：用整数位表示元素是否已选、资源是否拥有。风险点：同一位置不同掩码不能合并。

\textbf{位运算算术和低位技巧。}
代表案例：LeetCode 191 位 1 的个数、LeetCode 338 比特位计数、LeetCode 371 两整数之和。先看这些题的题面条件：题目要求不用普通算术或快速枚举子集。由此推出的处理方式是：用 lowbit、子集枚举或位加法模拟状态变化。风险点：移位优先级和有符号溢出需要谨慎。

\end{principlebox}

\subsection{案例矩阵}

本节案例覆盖 LeetCode 29 两数相除、LeetCode 50 Pow(x, n)、LeetCode 136 只出现一次的数字、LeetCode 137 只出现一次的数字 II、LeetCode 190 颠倒二进制位、LeetCode 191 位 1 的个数、LeetCode 201 数字范围按位与、LeetCode 231 2 的幂、LeetCode 260 只出现一次的数字 III、LeetCode 268 丢失的数字、LeetCode 287 寻找重复数、LeetCode 318 最大单词长度乘积、LeetCode 338 比特位计数、LeetCode 371 两整数之和、LeetCode 372 超级次方、LeetCode 393 UTF-8 编码验证、LeetCode 421 数组中两个数的最大异或值、LeetCode 1755 最接近目标值的子序列和。阅读时不要按题号背答案，而要先判断题面落在哪个分流场景：查询目标是什么，输入约束给了什么单调性或等价关系，暴力慢在哪里，优化结构保存了什么状态。

\begin{principlebox}{案例矩阵}
\textbf{LeetCode 29 两数相除。}
除法可以看成不断减去除数的倍增值。单点风险：除数为一、被除数为 int 最小值、结果溢出、符号处理。

\textbf{LeetCode 50 Pow(x, n)。}
指数 n 可以按二进制拆分，重复乘 x 的线性过程存在大量倍增结构。单点风险：n 为零、n 为负、int 最小值取反溢出、x 为零、浮点误差。

\textbf{LeetCode 136 只出现一次的数字。}
成对元素可以用异或抵消，剩下单个元素。单点风险：零、负数、唯一值在开头或结尾、所有其他值恰好两次。

\textbf{LeetCode 137 只出现一次的数字 II。}
除一个数字外，其余数字都出现三次，单纯异或不能抵消。单点风险：负数、符号位、目标为零、所有其他数恰好三次。

\textbf{LeetCode 190 颠倒二进制位。}
每一位的新位置由原位置镜像决定。单点风险：最高位、最低位、全零、全一、无符号语义。

\textbf{LeetCode 191 位 1 的个数。}
每次清掉最低位的一可以让循环次数等于一的个数。单点风险：零、只有最高位、一的个数很多、无符号输入。

\textbf{LeetCode 201 数字范围按位与。}
区间内所有数按位与后，只保留左右端公共二进制前缀。单点风险：left 等于 right、区间跨过二的幂边界、结果为零。

\textbf{LeetCode 231 2 的幂。}
二的幂在二进制中只有一个一。单点风险：n 为零、负数、int 最小值、最高位。

\textbf{LeetCode 260 只出现一次的数字 III。}
两个只出现一次的数异或后至少有一位不同。单点风险：两个答案一正一负、最低位分组、其他数恰好两次。

\textbf{LeetCode 268 丢失的数字。}
下标零到 n 与数组值相比，缺失值会在异或抵消后剩下。单点风险：缺失零、缺失 n、数组长度一、也可用求和但要防溢出。

\textbf{LeetCode 287 寻找重复数。}
数组值域使下标到值形成函数图，重复数是环入口。单点风险：重复数出现多次、不能排序、不能改数组、长度为 n 加一。

\textbf{LeetCode 318 最大单词长度乘积。}
两个单词没有公共字母时才可相乘，字母集合可压成位掩码。单点风险：重复字母、空字符串、多个单词同掩码、乘积范围。

\textbf{LeetCode 338 比特位计数。}
每个数的一的个数可由去掉最低位一后的较小数转移。单点风险：n 为零、二的幂、连续区间、表达式优先级。

\textbf{LeetCode 371 两整数之和。}
不用加号时，加法可以拆成无进位和与进位。单点风险：正负数、进位传播、有符号溢出、需要使用无符号中间值。

\textbf{LeetCode 372 超级次方。}
大指数以十进制数组给出，不能先转成普通整数。单点风险：指数为空、a 大于模数、指数含零、递归深度、模乘溢出。

\textbf{LeetCode 393 UTF-8 编码验证。}
每个字节的高位模式决定一个字符需要的后续字节数量。单点风险：非法首字节、后续字节不足、后续字节格式错误、单字节字符。

\textbf{LeetCode 421 数组中两个数的最大异或值。}
最大异或值由高位到低位尽量取一决定。单点风险：零、重复数、最高位、负数语义。

\textbf{LeetCode 1755 最接近目标值的子序列和。}
完整子序列选择是 \ensuremath{2^{n}}，拆成左右两半后各自枚举子集和。单点风险：负数、空子序列、目标为负、答案为零提前返回、两半大小。

\end{principlebox}

\subsection{共性落点}

\begin{principlebox}{本组共性落点}
\textbf{慢版本。} 暴力做法用集合、数组或递归保存状态，空间和比较成本较高。若状态本质是若干布尔选择，位掩码能把它压进整数。
\textbf{瓶颈。} 慢点在于状态复制和集合比较。位操作让包含、加入、删除、枚举子集变成常数或接近常数的机器指令。
\textbf{状态字段。} 掩码含义、当前位、目标位集合、状态转移和答案读取；负数或高位参与时要说明整数宽度。
\textbf{不变量。} 掩码的每一位必须有固定含义。状态去重时，位置相同但掩码不同不能合并，因为未来能力不同。
\textbf{实现检查。} 移位表达式加括号；使用无符号或足够宽类型；枚举子集时确认空集和全集是否包含。
\textbf{C++ 落地。} 使用无符号整数能减少移位歧义。位数超过 31 时用 \code{long long} 或 \code{std::bitset}。表达式加括号，避免位运算优先级误读。
\textbf{证明抓手。} 证明从逐位独立或子集归纳开始：被压缩到位里的信息足以决定未来转移。
\end{principlebox}

\subsection{案例推导}

\noindent\textbf{LeetCode 29 两数相除。}

\textbf{题目说明。} LeetCode 29 两数相除 的题面入口要先还原成具体任务：除法可以看成不断减去除数的倍增值。

\textbf{样例入口。} 拿 43 除以 5 手算，5 的倍增序列是 5、10、20、40，先减 40 商加 8，余 3 停止。

\textbf{暴力基线。} 慢版本先用集合、数组或布尔表保存状态，逐步执行题目动作。确认每个布尔字段的含义后，才能把它压缩成位，而不是直接写位运算公式。放到本题就是：先按“除法可以看成不断减去除数的倍增值”完整试慢版本；样例里已经能看到“规律是用二进制倍增找不超过被除数的最大除数倍数，符号和 int 最小值溢出要单独处理”，所以优化前必须先能说清慢版本枚举了哪些候选。

\textbf{暴力过程。} 拿 43 除以 5 手算，5 的倍增序列是 5、10、20、40，先减 40 商加 8，余 3 停止。

\textbf{重复工作。} 题面模型：除法可以看成不断减去除数的倍增值。暴力版的慢点要落到本题：慢版本会围绕“除法可以看成不断减去除数的倍增值”使用集合、计数或完整状态表。样例规律是“规律是用二进制倍增找不超过被除数的最大除数倍数，符号和 int 最小值溢出要单独处理”，说明哪些信息可以被逐位抵消、压缩或复用；位运算只是编码，不是证明本身。后面的优化只消掉这类重复工作，不能改变答案集合。

\textbf{优化条件。} 本题的关键观察是：把除数按二进制倍增，优先减去不超过被除数的最大倍数并累加商。这句话必须能翻译成可维护状态；如果题目条件改变后观察不再成立，就不能继续套原来的写法。

\textbf{相邻状态。} 规律是用二进制倍增找不超过被除数的最大除数倍数，符号和 int 最小值溢出要单独处理。把这个特例继续拆开看：每一位或每个掩码位都要对应题目里的一个布尔事实。位运算只是压缩表示；如果两个状态在位置相同但掩码不同，未来能力不同，就不能合并。

\textbf{状态复用。} 把集合状态压成掩码；包含、加入、删除、枚举子集都变成位操作，但每一位的题目含义不能丢。本题落点是：规律是用二进制倍增找不超过被除数的最大除数倍数，符号和 int 最小值溢出要单独处理。手算时只改被新输入影响的那一小部分，而不是回头重算完整候选。

\textbf{指针和字段。} 状态字段要能说清：掩码含义、当前位、目标位集合、状态转移和答案读取；负数或高位参与时要说明整数宽度。手算时按这个协议跟踪：拿 43 除以 5 手算，5 的倍增序列是 5、10、20、40，先减 40 商加 8，余 3 停止。然后按协议复查：手算时把整数写成二进制，并在每一位旁边标注含义。状态转移只允许改变题目动作允许改变的那些位。

\textbf{代码落点。} 移位表达式加括号；使用无符号或足够宽类型；枚举子集时确认空集和全集是否包含。关键代码行必须能逐句翻译成“旧状态怎样变成新状态”，否则就只是模板。

\textbf{正确性。} 证明从逐位独立或子集归纳开始：被压缩到位里的信息足以决定未来转移。

\textbf{边界实验。} 先用除数为一、被除数为 int 最小值、结果溢出、符号处理做反例检查。实验时覆盖除数为一、被除数为 int 最小值、结果溢出、符号处理，把数字写成二进制逐位检查。负数、最高位、空掩码和全集掩码都要单独测。

\textbf{复杂度变化。} 暴力用集合或数组保存状态可能需要更多空间和比较成本；位运算通常把每步状态更新压成常数或按位数计算，掩码 DP 仍按状态数增长。

\textbf{迁移追问。} 如果题目改成“它和快速幂类似，都利用二进制拆分减少重复加减”，先判断原状态是否仍然覆盖新答案集合，再决定是微调边界、改 value 语义，还是换算法。

\noindent\textbf{LeetCode 50 Pow(x, n)。}

\textbf{题目说明。} LeetCode 50 Pow(x, n) 的题面入口要先还原成具体任务：指数 n 可以按二进制拆分，重复乘 x 的线性过程存在大量倍增结构。

\textbf{样例入口。} 拿 x=2、n=10 手算，10 的二进制是 1010，所以答案是 \ensuremath{2^{8}} 乘 \ensuremath{2^{2}}。逐次乘 10 次太慢；每轮把底数平方、指数右移，当前位为 1 时乘入答案。

\textbf{暴力基线。} 慢版本先用集合、数组或布尔表保存状态，逐步执行题目动作。确认每个布尔字段的含义后，才能把它压缩成位，而不是直接写位运算公式。放到本题就是：先按“指数 n 可以按二进制拆分，重复乘 x 的线性过程存在大量倍增结构”完整试慢版本；样例里已经能看到“规律是指数按二进制拆分，负指数先转成倒数并用更宽整数处理”，所以优化前必须先能说清慢版本枚举了哪些候选。

\textbf{暴力过程。} 拿 x=2、n=10 手算，10 的二进制是 1010，所以答案是 \ensuremath{2^{8}} 乘 \ensuremath{2^{2}}。逐次乘 10 次太慢；每轮把底数平方、指数右移，当前位为 1 时乘入答案。

\textbf{重复工作。} 题面模型：指数 n 可以按二进制拆分，重复乘 x 的线性过程存在大量倍增结构。暴力版的慢点要落到本题：慢版本会围绕“指数 n 可以按二进制拆分，重复乘 x 的线性过程存在大量倍增结构”使用集合、计数或完整状态表。样例规律是“规律是指数按二进制拆分，负指数先转成倒数并用更宽整数处理”，说明哪些信息可以被逐位抵消、压缩或复用；位运算只是编码，不是证明本身。后面的优化只消掉这类重复工作，不能改变答案集合。

\textbf{优化条件。} 本题的关键观察是：快速幂每轮根据最低位决定是否乘入答案，同时底数平方、指数右移。这句话必须能翻译成可维护状态；如果题目条件改变后观察不再成立，就不能继续套原来的写法。

\textbf{相邻状态。} 规律是指数按二进制拆分，负指数先转成倒数并用更宽整数处理。把这个特例继续拆开看：每一位或每个掩码位都要对应题目里的一个布尔事实。位运算只是压缩表示；如果两个状态在位置相同但掩码不同，未来能力不同，就不能合并。

\textbf{状态复用。} 把集合状态压成掩码；包含、加入、删除、枚举子集都变成位操作，但每一位的题目含义不能丢。本题落点是：规律是指数按二进制拆分，负指数先转成倒数并用更宽整数处理。手算时只改被新输入影响的那一小部分，而不是回头重算完整候选。

\textbf{指针和字段。} 状态字段要能说清：掩码含义、当前位、目标位集合、状态转移和答案读取；负数或高位参与时要说明整数宽度。手算时按这个协议跟踪：拿 x=2、n=10 手算，10 的二进制是 1010，所以答案是 \ensuremath{2^{8}} 乘 \ensuremath{2^{2}}。逐次乘 10 次太慢；每轮把底数平方、指数右移，当前位为 1 时乘入答案。然后按协议复查：手算时把整数写成二进制，并在每一位旁边标注含义。状态转移只允许改变题目动作允许改变的那些位。

\textbf{代码落点。} 移位表达式加括号；使用无符号或足够宽类型；枚举子集时确认空集和全集是否包含。关键代码行必须能逐句翻译成“旧状态怎样变成新状态”，否则就只是模板。

\textbf{正确性。} 证明从逐位独立或子集归纳开始：被压缩到位里的信息足以决定未来转移。

\textbf{边界实验。} 先用n 为零、n 为负、int 最小值取反溢出、x 为零、浮点误差做反例检查。实验时覆盖n 为零、n 为负、int 最小值取反溢出、x 为零、浮点误差，把数字写成二进制逐位检查。负数、最高位、空掩码和全集掩码都要单独测。

\textbf{复杂度变化。} 暴力用集合或数组保存状态可能需要更多空间和比较成本；位运算通常把每步状态更新压成常数或按位数计算，掩码 DP 仍按状态数增长。

\textbf{迁移追问。} 如果题目改成“矩阵快速幂和模幂都是同一二进制拆分思想”，先判断原状态是否仍然覆盖新答案集合，再决定是微调边界、改 value 语义，还是换算法。

\noindent\textbf{LeetCode 136 只出现一次的数字。}

\textbf{题目说明。} LeetCode 136 只出现一次的数字 的题面入口要先还原成具体任务：成对元素可以用异或抵消，剩下单个元素。

\textbf{样例入口。} 拿 [4, 1, 2, 1, 2] 手算，1 和 1 异或抵消，2 和 2 抵消，最后剩 4。

\textbf{暴力基线。} 慢版本先用集合、数组或布尔表保存状态，逐步执行题目动作。确认每个布尔字段的含义后，才能把它压缩成位，而不是直接写位运算公式。放到本题就是：先按“成对元素可以用异或抵消，剩下单个元素”完整试慢版本；样例里已经能看到“规律是异或满足交换结合，且相同数异或后变成零，扫描顺序不影响结果；零和负数也只是按位参与同样运算”，所以优化前必须先能说清慢版本枚举了哪些候选。

\textbf{暴力过程。} 拿 [4, 1, 2, 1, 2] 手算，1 和 1 异或抵消，2 和 2 抵消，最后剩 4。

\textbf{重复工作。} 题面模型：成对元素可以用异或抵消，剩下单个元素。暴力版的慢点要落到本题：慢版本会围绕“成对元素可以用异或抵消，剩下单个元素”使用集合、计数或完整状态表。样例规律是“规律是异或满足交换结合，且相同数异或后变成零，扫描顺序不影响结果；零和负数也只是按位参与同样运算”，说明哪些信息可以被逐位抵消、压缩或复用；位运算只是编码，不是证明本身。后面的优化只消掉这类重复工作，不能改变答案集合。

\textbf{优化条件。} 本题的关键观察是：异或满足交换结合，扫描顺序不影响答案。这句话必须能翻译成可维护状态；如果题目条件改变后观察不再成立，就不能继续套原来的写法。

\textbf{相邻状态。} 规律是异或满足交换结合，且相同数异或后变成零，扫描顺序不影响结果；零和负数也只是按位参与同样运算。把这个特例继续拆开看：每一位或每个掩码位都要对应题目里的一个布尔事实。位运算只是压缩表示；如果两个状态在位置相同但掩码不同，未来能力不同，就不能合并。

\textbf{状态复用。} 把集合状态压成掩码；包含、加入、删除、枚举子集都变成位操作，但每一位的题目含义不能丢。本题落点是：规律是异或满足交换结合，且相同数异或后变成零，扫描顺序不影响结果；零和负数也只是按位参与同样运算。手算时只改被新输入影响的那一小部分，而不是回头重算完整候选。

\textbf{指针和字段。} 状态字段要能说清：掩码含义、当前位、目标位集合、状态转移和答案读取；负数或高位参与时要说明整数宽度。手算时按这个协议跟踪：拿 [4, 1, 2, 1, 2] 手算，1 和 1 异或抵消，2 和 2 抵消，最后剩 4。然后按协议复查：手算时把整数写成二进制，并在每一位旁边标注含义。状态转移只允许改变题目动作允许改变的那些位。

\textbf{代码落点。} 移位表达式加括号；使用无符号或足够宽类型；枚举子集时确认空集和全集是否包含。关键代码行必须能逐句翻译成“旧状态怎样变成新状态”，否则就只是模板。

\textbf{正确性。} 证明从逐位独立或子集归纳开始：被压缩到位里的信息足以决定未来转移。

\textbf{边界实验。} 先用零、负数、唯一值在开头或结尾、所有其他值恰好两次做反例检查。实验时覆盖零、负数、唯一值在开头或结尾、所有其他值恰好两次，把数字写成二进制逐位检查。负数、最高位、空掩码和全集掩码都要单独测。

\textbf{复杂度变化。} 暴力用集合或数组保存状态可能需要更多空间和比较成本；位运算通常把每步状态更新压成常数或按位数计算，掩码 DP 仍按状态数增长。

\textbf{迁移追问。} 如果题目改成“若其他值出现三次，要改为按位计数取模”，先判断原状态是否仍然覆盖新答案集合，再决定是微调边界、改 value 语义，还是换算法。

\noindent\textbf{LeetCode 137 只出现一次的数字 II。}

\textbf{题目说明。} LeetCode 137 只出现一次的数字 II 的题面入口要先还原成具体任务：除一个数字外，其余数字都出现三次，单纯异或不能抵消。

\textbf{样例入口。} 拿 [2, 2, 3, 2] 手算，每一位上 2 的贡献出现三次，对 3 取模后消失，只剩 3 的位。

\textbf{暴力基线。} 慢版本先用集合、数组或布尔表保存状态，逐步执行题目动作。确认每个布尔字段的含义后，才能把它压缩成位，而不是直接写位运算公式。放到本题就是：先按“除一个数字外，其余数字都出现三次，单纯异或不能抵消”完整试慢版本；样例里已经能看到“规律是逐位计数并对 3 取模，符号位也必须按固定整数宽度处理”，所以优化前必须先能说清慢版本枚举了哪些候选。

\textbf{暴力过程。} 拿 [2, 2, 3, 2] 手算，每一位上 2 的贡献出现三次，对 3 取模后消失，只剩 3 的位。

\textbf{重复工作。} 题面模型：除一个数字外，其余数字都出现三次，单纯异或不能抵消。暴力版的慢点要落到本题：慢版本会围绕“除一个数字外，其余数字都出现三次，单纯异或不能抵消”使用集合、计数或完整状态表。样例规律是“规律是逐位计数并对 3 取模，符号位也必须按固定整数宽度处理”，说明哪些信息可以被逐位抵消、压缩或复用；位运算只是编码，不是证明本身。后面的优化只消掉这类重复工作，不能改变答案集合。

\textbf{优化条件。} 本题的关键观察是：逐位统计一的数量，对三取模后还原目标数字。这句话必须能翻译成可维护状态；如果题目条件改变后观察不再成立，就不能继续套原来的写法。

\textbf{相邻状态。} 规律是逐位计数并对 3 取模，符号位也必须按固定整数宽度处理。把这个特例继续拆开看：每一位或每个掩码位都要对应题目里的一个布尔事实。位运算只是压缩表示；如果两个状态在位置相同但掩码不同，未来能力不同，就不能合并。

\textbf{状态复用。} 把集合状态压成掩码；包含、加入、删除、枚举子集都变成位操作，但每一位的题目含义不能丢。本题落点是：规律是逐位计数并对 3 取模，符号位也必须按固定整数宽度处理。手算时只改被新输入影响的那一小部分，而不是回头重算完整候选。

\textbf{指针和字段。} 状态字段要能说清：掩码含义、当前位、目标位集合、状态转移和答案读取；负数或高位参与时要说明整数宽度。手算时按这个协议跟踪：拿 [2, 2, 3, 2] 手算，每一位上 2 的贡献出现三次，对 3 取模后消失，只剩 3 的位。然后按协议复查：手算时把整数写成二进制，并在每一位旁边标注含义。状态转移只允许改变题目动作允许改变的那些位。

\textbf{代码落点。} 移位表达式加括号；使用无符号或足够宽类型；枚举子集时确认空集和全集是否包含。关键代码行必须能逐句翻译成“旧状态怎样变成新状态”，否则就只是模板。

\textbf{正确性。} 证明从逐位独立或子集归纳开始：被压缩到位里的信息足以决定未来转移。

\textbf{边界实验。} 先用负数、符号位、目标为零、所有其他数恰好三次做反例检查。实验时覆盖负数、符号位、目标为零、所有其他数恰好三次，把数字写成二进制逐位检查。负数、最高位、空掩码和全集掩码都要单独测。

\textbf{复杂度变化。} 暴力用集合或数组保存状态可能需要更多空间和比较成本；位运算通常把每步状态更新压成常数或按位数计算，掩码 DP 仍按状态数增长。

\textbf{迁移追问。} 如果题目改成“若其他数出现 k 次，思路仍是按位计数取模”，先判断原状态是否仍然覆盖新答案集合，再决定是微调边界、改 value 语义，还是换算法。

\noindent\textbf{LeetCode 190 颠倒二进制位。}

\textbf{题目说明。} LeetCode 190 颠倒二进制位 的题面入口要先还原成具体任务：每一位的新位置由原位置镜像决定。

\textbf{样例入口。} 拿低位是 101 的数手算，颠倒时每次把结果左移一位，再接上当前最低位；原数右移继续取下一位。

\textbf{暴力基线。} 慢版本先用集合、数组或布尔表保存状态，逐步执行题目动作。确认每个布尔字段的含义后，才能把它压缩成位，而不是直接写位运算公式。放到本题就是：先按“每一位的新位置由原位置镜像决定”完整试慢版本；样例里已经能看到“规律是循环 32 次处理固定宽度，不能因为高位暂时为 0 就提前结束”，所以优化前必须先能说清慢版本枚举了哪些候选。

\textbf{暴力过程。} 拿低位是 101 的数手算，颠倒时每次把结果左移一位，再接上当前最低位；原数右移继续取下一位。

\textbf{重复工作。} 题面模型：每一位的新位置由原位置镜像决定。暴力版的慢点要落到本题：慢版本会围绕“每一位的新位置由原位置镜像决定”使用集合、计数或完整状态表。样例规律是“规律是循环 32 次处理固定宽度，不能因为高位暂时为 0 就提前结束”，说明哪些信息可以被逐位抵消、压缩或复用；位运算只是编码，不是证明本身。后面的优化只消掉这类重复工作，不能改变答案集合。

\textbf{优化条件。} 本题的关键观察是：循环三十二次，把结果左移并接上当前最低位，再右移输入。这句话必须能翻译成可维护状态；如果题目条件改变后观察不再成立，就不能继续套原来的写法。

\textbf{相邻状态。} 规律是循环 32 次处理固定宽度，不能因为高位暂时为 0 就提前结束。把这个特例继续拆开看：每一位或每个掩码位都要对应题目里的一个布尔事实。位运算只是压缩表示；如果两个状态在位置相同但掩码不同，未来能力不同，就不能合并。

\textbf{状态复用。} 把集合状态压成掩码；包含、加入、删除、枚举子集都变成位操作，但每一位的题目含义不能丢。本题落点是：规律是循环 32 次处理固定宽度，不能因为高位暂时为 0 就提前结束。手算时只改被新输入影响的那一小部分，而不是回头重算完整候选。

\textbf{指针和字段。} 状态字段要能说清：掩码含义、当前位、目标位集合、状态转移和答案读取；负数或高位参与时要说明整数宽度。手算时按这个协议跟踪：拿低位是 101 的数手算，颠倒时每次把结果左移一位，再接上当前最低位；原数右移继续取下一位。然后按协议复查：手算时把整数写成二进制，并在每一位旁边标注含义。状态转移只允许改变题目动作允许改变的那些位。

\textbf{代码落点。} 移位表达式加括号；使用无符号或足够宽类型；枚举子集时确认空集和全集是否包含。关键代码行必须能逐句翻译成“旧状态怎样变成新状态”，否则就只是模板。

\textbf{正确性。} 证明从逐位独立或子集归纳开始：被压缩到位里的信息足以决定未来转移。

\textbf{边界实验。} 先用最高位、最低位、全零、全一、无符号语义做反例检查。实验时覆盖最高位、最低位、全零、全一、无符号语义，把数字写成二进制逐位检查。负数、最高位、空掩码和全集掩码都要单独测。

\textbf{复杂度变化。} 暴力用集合或数组保存状态可能需要更多空间和比较成本；位运算通常把每步状态更新压成常数或按位数计算，掩码 DP 仍按状态数增长。

\textbf{迁移追问。} 如果题目改成“若多次调用，可以按字节预处理反转表”，先判断原状态是否仍然覆盖新答案集合，再决定是微调边界、改 value 语义，还是换算法。

\noindent\textbf{LeetCode 191 位 1 的个数。}

\textbf{题目说明。} LeetCode 191 位 1 的个数 的题面入口要先还原成具体任务：每次清掉最低位的一可以让循环次数等于一的个数。

\textbf{样例入口。} 拿 11 的二进制 1011 手算，n\&(n-1) 第一次删掉最低位 1 得 1010，第二次得 1000，第三次得 0。

\textbf{暴力基线。} 慢版本先用集合、数组或布尔表保存状态，逐步执行题目动作。确认每个布尔字段的含义后，才能把它压缩成位，而不是直接写位运算公式。放到本题就是：先按“每次清掉最低位的一可以让循环次数等于一的个数”完整试慢版本；样例里已经能看到“规律是每次操作恰好删除一个 1，所以循环次数就是 1 的个数”，所以优化前必须先能说清慢版本枚举了哪些候选。

\textbf{暴力过程。} 拿 11 的二进制 1011 手算，n\&(n-1) 第一次删掉最低位 1 得 1010，第二次得 1000，第三次得 0。

\textbf{重复工作。} 题面模型：每次清掉最低位的一可以让循环次数等于一的个数。暴力版的慢点要落到本题：慢版本会围绕“每次清掉最低位的一可以让循环次数等于一的个数”使用集合、计数或完整状态表。样例规律是“规律是每次操作恰好删除一个 1，所以循环次数就是 1 的个数”，说明哪些信息可以被逐位抵消、压缩或复用；位运算只是编码，不是证明本身。后面的优化只消掉这类重复工作，不能改变答案集合。

\textbf{优化条件。} 本题的关键观察是：使用 n \& (n - 1) 删除最低位一，比逐位检查更直接。这句话必须能翻译成可维护状态；如果题目条件改变后观察不再成立，就不能继续套原来的写法。

\textbf{相邻状态。} 规律是每次操作恰好删除一个 1，所以循环次数就是 1 的个数。把这个特例继续拆开看：每一位或每个掩码位都要对应题目里的一个布尔事实。位运算只是压缩表示；如果两个状态在位置相同但掩码不同，未来能力不同，就不能合并。

\textbf{状态复用。} 把集合状态压成掩码；包含、加入、删除、枚举子集都变成位操作，但每一位的题目含义不能丢。本题落点是：规律是每次操作恰好删除一个 1，所以循环次数就是 1 的个数。手算时只改被新输入影响的那一小部分，而不是回头重算完整候选。

\textbf{指针和字段。} 状态字段要能说清：掩码含义、当前位、目标位集合、状态转移和答案读取；负数或高位参与时要说明整数宽度。手算时按这个协议跟踪：拿 11 的二进制 1011 手算，n\&(n-1) 第一次删掉最低位 1 得 1010，第二次得 1000，第三次得 0。然后按协议复查：手算时把整数写成二进制，并在每一位旁边标注含义。状态转移只允许改变题目动作允许改变的那些位。

\textbf{代码落点。} 移位表达式加括号；使用无符号或足够宽类型；枚举子集时确认空集和全集是否包含。关键代码行必须能逐句翻译成“旧状态怎样变成新状态”，否则就只是模板。

\textbf{正确性。} 证明从逐位独立或子集归纳开始：被压缩到位里的信息足以决定未来转移。

\textbf{边界实验。} 先用零、只有最高位、一的个数很多、无符号输入做反例检查。实验时覆盖零、只有最高位、一的个数很多、无符号输入，把数字写成二进制逐位检查。负数、最高位、空掩码和全集掩码都要单独测。

\textbf{复杂度变化。} 暴力用集合或数组保存状态可能需要更多空间和比较成本；位运算通常把每步状态更新压成常数或按位数计算，掩码 DP 仍按状态数增长。

\textbf{迁移追问。} 如果题目改成“比特位计数可以把这个过程 DP 化到一整段数”，先判断原状态是否仍然覆盖新答案集合，再决定是微调边界、改 value 语义，还是换算法。

\noindent\textbf{LeetCode 201 数字范围按位与。}

\textbf{题目说明。} LeetCode 201 数字范围按位与 的题面入口要先还原成具体任务：区间内所有数按位与后，只保留左右端公共二进制前缀。

\textbf{样例入口。} 拿区间 [5, 7]，二进制是 101、110、111，只有最高公共前缀 1 能留下，低位在区间内会出现 0 和 1，被按位与清零。

\textbf{暴力基线。} 慢版本先用集合、数组或布尔表保存状态，逐步执行题目动作。确认每个布尔字段的含义后，才能把它压缩成位，而不是直接写位运算公式。放到本题就是：先按“区间内所有数按位与后，只保留左右端公共二进制前缀”完整试慢版本；样例里已经能看到“规律是不断右移 left 和 right 找公共前缀，再移回原位”，所以优化前必须先能说清慢版本枚举了哪些候选。

\textbf{暴力过程。} 拿区间 [5, 7]，二进制是 101、110、111，只有最高公共前缀 1 能留下，低位在区间内会出现 0 和 1，被按位与清零。

\textbf{重复工作。} 题面模型：区间内所有数按位与后，只保留左右端公共二进制前缀。暴力版的慢点要落到本题：慢版本会围绕“区间内所有数按位与后，只保留左右端公共二进制前缀”使用集合、计数或完整状态表。样例规律是“规律是不断右移 left 和 right 找公共前缀，再移回原位”，说明哪些信息可以被逐位抵消、压缩或复用；位运算只是编码，不是证明本身。后面的优化只消掉这类重复工作，不能改变答案集合。

\textbf{优化条件。} 本题的关键观察是：不断右移 left 和 right，直到二者相等，再把公共前缀移回原位。这句话必须能翻译成可维护状态；如果题目条件改变后观察不再成立，就不能继续套原来的写法。

\textbf{相邻状态。} 规律是不断右移 left 和 right 找公共前缀，再移回原位。把这个特例继续拆开看：每一位或每个掩码位都要对应题目里的一个布尔事实。位运算只是压缩表示；如果两个状态在位置相同但掩码不同，未来能力不同，就不能合并。

\textbf{状态复用。} 把集合状态压成掩码；包含、加入、删除、枚举子集都变成位操作，但每一位的题目含义不能丢。本题落点是：规律是不断右移 left 和 right 找公共前缀，再移回原位。手算时只改被新输入影响的那一小部分，而不是回头重算完整候选。

\textbf{指针和字段。} 状态字段要能说清：掩码含义、当前位、目标位集合、状态转移和答案读取；负数或高位参与时要说明整数宽度。手算时按这个协议跟踪：拿区间 [5, 7]，二进制是 101、110、111，只有最高公共前缀 1 能留下，低位在区间内会出现 0 和 1，被按位与清零。然后按协议复查：手算时把整数写成二进制，并在每一位旁边标注含义。状态转移只允许改变题目动作允许改变的那些位。

\textbf{代码落点。} 移位表达式加括号；使用无符号或足够宽类型；枚举子集时确认空集和全集是否包含。关键代码行必须能逐句翻译成“旧状态怎样变成新状态”，否则就只是模板。

\textbf{正确性。} 证明从逐位独立或子集归纳开始：被压缩到位里的信息足以决定未来转移。

\textbf{边界实验。} 先用left 等于 right、区间跨过二的幂边界、结果为零做反例检查。实验时覆盖left 等于 right、区间跨过二的幂边界、结果为零，把数字写成二进制逐位检查。负数、最高位、空掩码和全集掩码都要单独测。

\textbf{复杂度变化。} 暴力用集合或数组保存状态可能需要更多空间和比较成本；位运算通常把每步状态更新压成常数或按位数计算，掩码 DP 仍按状态数增长。

\textbf{迁移追问。} 如果题目改成“也可以不断清除 right 的最低位一，直到 right 不大于 left”，先判断原状态是否仍然覆盖新答案集合，再决定是微调边界、改 value 语义，还是换算法。

\noindent\textbf{LeetCode 231 2 的幂。}

\textbf{题目说明。} LeetCode 231 2 的幂 的题面入口要先还原成具体任务：二的幂在二进制中只有一个一。

\textbf{样例入口。} 拿 8 的二进制 1000 手算，8\&(7)=1000\&0111=0；拿 6 的 110，6\&5=100 不为 0。

\textbf{暴力基线。} 慢版本先用集合、数组或布尔表保存状态，逐步执行题目动作。确认每个布尔字段的含义后，才能把它压缩成位，而不是直接写位运算公式。放到本题就是：先按“二的幂在二进制中只有一个一”完整试慢版本；样例里已经能看到“规律是 2 的幂只有一个 1，且 n 必须为正；0 和负数不能只靠按位与判断”，所以优化前必须先能说清慢版本枚举了哪些候选。

\textbf{暴力过程。} 拿 8 的二进制 1000 手算，8\&(7)=1000\&0111=0；拿 6 的 110，6\&5=100 不为 0。

\textbf{重复工作。} 题面模型：二的幂在二进制中只有一个一。暴力版的慢点要落到本题：慢版本会围绕“二的幂在二进制中只有一个一”使用集合、计数或完整状态表。样例规律是“规律是 2 的幂只有一个 1，且 n 必须为正；0 和负数不能只靠按位与判断”，说明哪些信息可以被逐位抵消、压缩或复用；位运算只是编码，不是证明本身。后面的优化只消掉这类重复工作，不能改变答案集合。

\textbf{优化条件。} 本题的关键观察是：正数 n 满足 n 与 n 减一按位与为零。这句话必须能翻译成可维护状态；如果题目条件改变后观察不再成立，就不能继续套原来的写法。

\textbf{相邻状态。} 规律是 2 的幂只有一个 1，且 n 必须为正；0 和负数不能只靠按位与判断。把这个特例继续拆开看：每一位或每个掩码位都要对应题目里的一个布尔事实。位运算只是压缩表示；如果两个状态在位置相同但掩码不同，未来能力不同，就不能合并。

\textbf{状态复用。} 把集合状态压成掩码；包含、加入、删除、枚举子集都变成位操作，但每一位的题目含义不能丢。本题落点是：规律是 2 的幂只有一个 1，且 n 必须为正；0 和负数不能只靠按位与判断。手算时只改被新输入影响的那一小部分，而不是回头重算完整候选。

\textbf{指针和字段。} 状态字段要能说清：掩码含义、当前位、目标位集合、状态转移和答案读取；负数或高位参与时要说明整数宽度。手算时按这个协议跟踪：拿 8 的二进制 1000 手算，8\&(7)=1000\&0111=0；拿 6 的 110，6\&5=100 不为 0。然后按协议复查：手算时把整数写成二进制，并在每一位旁边标注含义。状态转移只允许改变题目动作允许改变的那些位。

\textbf{代码落点。} 移位表达式加括号；使用无符号或足够宽类型；枚举子集时确认空集和全集是否包含。关键代码行必须能逐句翻译成“旧状态怎样变成新状态”，否则就只是模板。

\textbf{正确性。} 证明从逐位独立或子集归纳开始：被压缩到位里的信息足以决定未来转移。

\textbf{边界实验。} 先用n 为零、负数、int 最小值、最高位做反例检查。实验时覆盖n 为零、负数、int 最小值、最高位，把数字写成二进制逐位检查。负数、最高位、空掩码和全集掩码都要单独测。

\textbf{复杂度变化。} 暴力用集合或数组保存状态可能需要更多空间和比较成本；位运算通常把每步状态更新压成常数或按位数计算，掩码 DP 仍按状态数增长。

\textbf{迁移追问。} 如果题目改成“判断四的幂还要额外限制一所在的位置为偶数位”，先判断原状态是否仍然覆盖新答案集合，再决定是微调边界、改 value 语义，还是换算法。

\noindent\textbf{LeetCode 260 只出现一次的数字 III。}

\textbf{题目说明。} LeetCode 260 只出现一次的数字 III 的题面入口要先还原成具体任务：两个只出现一次的数异或后至少有一位不同。

\textbf{样例入口。} 拿 [1, 2, 1, 3, 2, 5] 手算，全部异或后只剩 3 和 5 的异或结果，取最低不同位可把 3 和 5 分到两组；每组内部重复数再异或抵消。

\textbf{暴力基线。} 慢版本先用集合、数组或布尔表保存状态，逐步执行题目动作。确认每个布尔字段的含义后，才能把它压缩成位，而不是直接写位运算公式。放到本题就是：先按“两个只出现一次的数异或后至少有一位不同”完整试慢版本；样例里已经能看到“规律是两个答案至少有一位不同，用这位分组后退化成两个只出现一次的问题”，所以优化前必须先能说清慢版本枚举了哪些候选。

\textbf{暴力过程。} 拿 [1, 2, 1, 3, 2, 5] 手算，全部异或后只剩 3 和 5 的异或结果，取最低不同位可把 3 和 5 分到两组；每组内部重复数再异或抵消。

\textbf{重复工作。} 题面模型：两个只出现一次的数异或后至少有一位不同。暴力版的慢点要落到本题：慢版本会围绕“两个只出现一次的数异或后至少有一位不同”使用集合、计数或完整状态表。样例规律是“规律是两个答案至少有一位不同，用这位分组后退化成两个只出现一次的问题”，说明哪些信息可以被逐位抵消、压缩或复用；位运算只是编码，不是证明本身。后面的优化只消掉这类重复工作，不能改变答案集合。

\textbf{优化条件。} 本题的关键观察是：用最低不同位把数组分成两组，每组内部再异或抵消。这句话必须能翻译成可维护状态；如果题目条件改变后观察不再成立，就不能继续套原来的写法。

\textbf{相邻状态。} 规律是两个答案至少有一位不同，用这位分组后退化成两个只出现一次的问题。把这个特例继续拆开看：每一位或每个掩码位都要对应题目里的一个布尔事实。位运算只是压缩表示；如果两个状态在位置相同但掩码不同，未来能力不同，就不能合并。

\textbf{状态复用。} 把集合状态压成掩码；包含、加入、删除、枚举子集都变成位操作，但每一位的题目含义不能丢。本题落点是：规律是两个答案至少有一位不同，用这位分组后退化成两个只出现一次的问题。手算时只改被新输入影响的那一小部分，而不是回头重算完整候选。

\textbf{指针和字段。} 状态字段要能说清：掩码含义、当前位、目标位集合、状态转移和答案读取；负数或高位参与时要说明整数宽度。手算时按这个协议跟踪：拿 [1, 2, 1, 3, 2, 5] 手算，全部异或后只剩 3 和 5 的异或结果，取最低不同位可把 3 和 5 分到两组；每组内部重复数再异或抵消。然后按协议复查：手算时把整数写成二进制，并在每一位旁边标注含义。状态转移只允许改变题目动作允许改变的那些位。

\textbf{代码落点。} 移位表达式加括号；使用无符号或足够宽类型；枚举子集时确认空集和全集是否包含。关键代码行必须能逐句翻译成“旧状态怎样变成新状态”，否则就只是模板。

\textbf{正确性。} 证明从逐位独立或子集归纳开始：被压缩到位里的信息足以决定未来转移。

\textbf{边界实验。} 先用两个答案一正一负、最低位分组、其他数恰好两次做反例检查。实验时覆盖两个答案一正一负、最低位分组、其他数恰好两次，把数字写成二进制逐位检查。负数、最高位、空掩码和全集掩码都要单独测。

\textbf{复杂度变化。} 暴力用集合或数组保存状态可能需要更多空间和比较成本；位运算通常把每步状态更新压成常数或按位数计算，掩码 DP 仍按状态数增长。

\textbf{迁移追问。} 如果题目改成“它是在异或抵消基础上增加分组位”，先判断原状态是否仍然覆盖新答案集合，再决定是微调边界、改 value 语义，还是换算法。

\noindent\textbf{LeetCode 268 丢失的数字。}

\textbf{题目说明。} LeetCode 268 丢失的数字 的题面入口要先还原成具体任务：下标零到 n 与数组值相比，缺失值会在异或抵消后剩下。

\textbf{样例入口。} 拿 [3, 0, 1] 手算，0、1、3 与完整下标 0、1、2、3 一起异或，重复的 0、1、3 抵消，剩 2。

\textbf{暴力基线。} 慢版本先用集合、数组或布尔表保存状态，逐步执行题目动作。确认每个布尔字段的含义后，才能把它压缩成位，而不是直接写位运算公式。放到本题就是：先按“下标零到 n 与数组值相比，缺失值会在异或抵消后剩下”完整试慢版本；样例里已经能看到“规律是缺失数可以由异或抵消得到，也可用求和但要注意溢出”，所以优化前必须先能说清慢版本枚举了哪些候选。

\textbf{暴力过程。} 拿 [3, 0, 1] 手算，0、1、3 与完整下标 0、1、2、3 一起异或，重复的 0、1、3 抵消，剩 2。

\textbf{重复工作。} 题面模型：下标零到 n 与数组值相比，缺失值会在异或抵消后剩下。暴力版的慢点要落到本题：慢版本会围绕“下标零到 n 与数组值相比，缺失值会在异或抵消后剩下”使用集合、计数或完整状态表。样例规律是“规律是缺失数可以由异或抵消得到，也可用求和但要注意溢出”，说明哪些信息可以被逐位抵消、压缩或复用；位运算只是编码，不是证明本身。后面的优化只消掉这类重复工作，不能改变答案集合。

\textbf{优化条件。} 本题的关键观察是：把所有下标和所有值一起异或，重复出现的数抵消。这句话必须能翻译成可维护状态；如果题目条件改变后观察不再成立，就不能继续套原来的写法。

\textbf{相邻状态。} 规律是缺失数可以由异或抵消得到，也可用求和但要注意溢出。把这个特例继续拆开看：每一位或每个掩码位都要对应题目里的一个布尔事实。位运算只是压缩表示；如果两个状态在位置相同但掩码不同，未来能力不同，就不能合并。

\textbf{状态复用。} 把集合状态压成掩码；包含、加入、删除、枚举子集都变成位操作，但每一位的题目含义不能丢。本题落点是：规律是缺失数可以由异或抵消得到，也可用求和但要注意溢出。手算时只改被新输入影响的那一小部分，而不是回头重算完整候选。

\textbf{指针和字段。} 状态字段要能说清：掩码含义、当前位、目标位集合、状态转移和答案读取；负数或高位参与时要说明整数宽度。手算时按这个协议跟踪：拿 [3, 0, 1] 手算，0、1、3 与完整下标 0、1、2、3 一起异或，重复的 0、1、3 抵消，剩 2。然后按协议复查：手算时把整数写成二进制，并在每一位旁边标注含义。状态转移只允许改变题目动作允许改变的那些位。

\textbf{代码落点。} 移位表达式加括号；使用无符号或足够宽类型；枚举子集时确认空集和全集是否包含。关键代码行必须能逐句翻译成“旧状态怎样变成新状态”，否则就只是模板。

\textbf{正确性。} 证明从逐位独立或子集归纳开始：被压缩到位里的信息足以决定未来转移。

\textbf{边界实验。} 先用缺失零、缺失 n、数组长度一、也可用求和但要防溢出做反例检查。实验时覆盖缺失零、缺失 n、数组长度一、也可用求和但要防溢出，把数字写成二进制逐位检查。负数、最高位、空掩码和全集掩码都要单独测。

\textbf{复杂度变化。} 暴力用集合或数组保存状态可能需要更多空间和比较成本；位运算通常把每步状态更新压成常数或按位数计算，掩码 DP 仍按状态数增长。

\textbf{迁移追问。} 如果题目改成“异或法适合展示值域完整且只有一个缺口的结构”，先判断原状态是否仍然覆盖新答案集合，再决定是微调边界、改 value 语义，还是换算法。

\noindent\textbf{LeetCode 287 寻找重复数。}

\textbf{题目说明。} LeetCode 287 寻找重复数 的题面入口要先还原成具体任务：数组值域使下标到值形成函数图，重复数是环入口。

\textbf{样例入口。} 拿 [1, 3, 4, 2, 2] 手算，把值当成 next 指针会形成 1->3->2->4->2 的环，入环点 2 就是重复数。

\textbf{暴力基线。} 慢版本先用集合、数组或布尔表保存状态，逐步执行题目动作。确认每个布尔字段的含义后，才能把它压缩成位，而不是直接写位运算公式。放到本题就是：先按“数组值域使下标到值形成函数图，重复数是环入口”完整试慢版本；样例里已经能看到“规律是长度 n+1、值域 1..n 强制形成环；不能修改数组时可用快慢指针或值域二分”，所以优化前必须先能说清慢版本枚举了哪些候选。

\textbf{暴力过程。} 拿 [1, 3, 4, 2, 2] 手算，把值当成 next 指针会形成 1->3->2->4->2 的环，入环点 2 就是重复数。

\textbf{重复工作。} 题面模型：数组值域使下标到值形成函数图，重复数是环入口。暴力版的慢点要落到本题：慢版本会围绕“数组值域使下标到值形成函数图，重复数是环入口”使用集合、计数或完整状态表。样例规律是“规律是长度 n+1、值域 1..n 强制形成环；不能修改数组时可用快慢指针或值域二分”，说明哪些信息可以被逐位抵消、压缩或复用；位运算只是编码，不是证明本身。后面的优化只消掉这类重复工作，不能改变答案集合。

\textbf{优化条件。} 本题的关键观察是：快慢指针不修改数组且常数空间；值域二分用抽屉原理。这句话必须能翻译成可维护状态；如果题目条件改变后观察不再成立，就不能继续套原来的写法。

\textbf{相邻状态。} 规律是长度 n+1、值域 1..n 强制形成环；不能修改数组时可用快慢指针或值域二分。把这个特例继续拆开看：每一位或每个掩码位都要对应题目里的一个布尔事实。位运算只是压缩表示；如果两个状态在位置相同但掩码不同，未来能力不同，就不能合并。

\textbf{状态复用。} 把集合状态压成掩码；包含、加入、删除、枚举子集都变成位操作，但每一位的题目含义不能丢。本题落点是：规律是长度 n+1、值域 1..n 强制形成环；不能修改数组时可用快慢指针或值域二分。手算时只改被新输入影响的那一小部分，而不是回头重算完整候选。

\textbf{指针和字段。} 状态字段要能说清：掩码含义、当前位、目标位集合、状态转移和答案读取；负数或高位参与时要说明整数宽度。手算时按这个协议跟踪：拿 [1, 3, 4, 2, 2] 手算，把值当成 next 指针会形成 1->3->2->4->2 的环，入环点 2 就是重复数。然后按协议复查：手算时把整数写成二进制，并在每一位旁边标注含义。状态转移只允许改变题目动作允许改变的那些位。

\textbf{代码落点。} 移位表达式加括号；使用无符号或足够宽类型；枚举子集时确认空集和全集是否包含。关键代码行必须能逐句翻译成“旧状态怎样变成新状态”，否则就只是模板。

\textbf{正确性。} 证明从逐位独立或子集归纳开始：被压缩到位里的信息足以决定未来转移。

\textbf{边界实验。} 先用重复数出现多次、不能排序、不能改数组、长度为 n 加一做反例检查。实验时覆盖重复数出现多次、不能排序、不能改数组、长度为 n 加一，把数字写成二进制逐位检查。负数、最高位、空掩码和全集掩码都要单独测。

\textbf{复杂度变化。} 暴力用集合或数组保存状态可能需要更多空间和比较成本；位运算通常把每步状态更新压成常数或按位数计算，掩码 DP 仍按状态数增长。

\textbf{迁移追问。} 如果题目改成“快慢指针要求值都能作为合法下标”，先判断原状态是否仍然覆盖新答案集合，再决定是微调边界、改 value 语义，还是换算法。

\noindent\textbf{LeetCode 318 最大单词长度乘积。}

\textbf{题目说明。} LeetCode 318 最大单词长度乘积 的题面入口要先还原成具体任务：两个单词没有公共字母时才可相乘，字母集合可压成位掩码。

\textbf{样例入口。} 拿 abcw 和 xtfn 手算，两个单词的 26 位掩码按位与为 0，说明没有公共字符，乘积 4*4 可作为候选。

\textbf{暴力基线。} 慢版本先用集合、数组或布尔表保存状态，逐步执行题目动作。确认每个布尔字段的含义后，才能把它压缩成位，而不是直接写位运算公式。放到本题就是：先按“两个单词没有公共字母时才可相乘，字母集合可压成位掩码”完整试慢版本；样例里已经能看到“规律是每个单词压成字符集合掩码，重复字母不增加位，掩码相交才排除”，所以优化前必须先能说清慢版本枚举了哪些候选。

\textbf{暴力过程。} 拿 abcw 和 xtfn 手算，两个单词的 26 位掩码按位与为 0，说明没有公共字符，乘积 4*4 可作为候选。

\textbf{重复工作。} 题面模型：两个单词没有公共字母时才可相乘，字母集合可压成位掩码。暴力版的慢点要落到本题：慢版本会围绕“两个单词没有公共字母时才可相乘，字母集合可压成位掩码”使用集合、计数或完整状态表。样例规律是“规律是每个单词压成字符集合掩码，重复字母不增加位，掩码相交才排除”，说明哪些信息可以被逐位抵消、压缩或复用；位运算只是编码，不是证明本身。后面的优化只消掉这类重复工作，不能改变答案集合。

\textbf{优化条件。} 本题的关键观察是：每个单词生成 26 位掩码，两个掩码按位与为零表示无交集。这句话必须能翻译成可维护状态；如果题目条件改变后观察不再成立，就不能继续套原来的写法。

\textbf{相邻状态。} 规律是每个单词压成字符集合掩码，重复字母不增加位，掩码相交才排除。把这个特例继续拆开看：每一位或每个掩码位都要对应题目里的一个布尔事实。位运算只是压缩表示；如果两个状态在位置相同但掩码不同，未来能力不同，就不能合并。

\textbf{状态复用。} 把集合状态压成掩码；包含、加入、删除、枚举子集都变成位操作，但每一位的题目含义不能丢。本题落点是：规律是每个单词压成字符集合掩码，重复字母不增加位，掩码相交才排除。手算时只改被新输入影响的那一小部分，而不是回头重算完整候选。

\textbf{指针和字段。} 状态字段要能说清：掩码含义、当前位、目标位集合、状态转移和答案读取；负数或高位参与时要说明整数宽度。手算时按这个协议跟踪：拿 abcw 和 xtfn 手算，两个单词的 26 位掩码按位与为 0，说明没有公共字符，乘积 4*4 可作为候选。然后按协议复查：手算时把整数写成二进制，并在每一位旁边标注含义。状态转移只允许改变题目动作允许改变的那些位。

\textbf{代码落点。} 移位表达式加括号；使用无符号或足够宽类型；枚举子集时确认空集和全集是否包含。关键代码行必须能逐句翻译成“旧状态怎样变成新状态”，否则就只是模板。

\textbf{正确性。} 证明从逐位独立或子集归纳开始：被压缩到位里的信息足以决定未来转移。

\textbf{边界实验。} 先用重复字母、空字符串、多个单词同掩码、乘积范围做反例检查。实验时覆盖重复字母、空字符串、多个单词同掩码、乘积范围，把数字写成二进制逐位检查。负数、最高位、空掩码和全集掩码都要单独测。

\textbf{复杂度变化。} 暴力用集合或数组保存状态可能需要更多空间和比较成本；位运算通常把每步状态更新压成常数或按位数计算，掩码 DP 仍按状态数增长。

\textbf{迁移追问。} 如果题目改成“若字符集大于整数位数，需要 bitset 或哈希集合”，先判断原状态是否仍然覆盖新答案集合，再决定是微调边界、改 value 语义，还是换算法。

\noindent\textbf{LeetCode 338 比特位计数。}

\textbf{题目说明。} LeetCode 338 比特位计数 的题面入口要先还原成具体任务：每个数的一的个数可由去掉最低位一后的较小数转移。

\textbf{样例入口。} 拿 n=5 手算，5 的二进制 101，比 4 多一个最低位 1，所以 bits[5]=bits[4]+1；也可用 5\&(4)=4 得到 bits[5]=bits[4]+1。

\textbf{暴力基线。} 慢版本先用集合、数组或布尔表保存状态，逐步执行题目动作。确认每个布尔字段的含义后，才能把它压缩成位，而不是直接写位运算公式。放到本题就是：先按“每个数的一的个数可由去掉最低位一后的较小数转移”完整试慢版本；样例里已经能看到“规律是每个数的答案来自去掉最低位 1 或右移一位后的更小数”，所以优化前必须先能说清慢版本枚举了哪些候选。

\textbf{暴力过程。} 拿 n=5 手算，5 的二进制 101，比 4 多一个最低位 1，所以 bits[5]=bits[4]+1；也可用 5\&(4)=4 得到 bits[5]=bits[4]+1。

\textbf{重复工作。} 题面模型：每个数的一的个数可由去掉最低位一后的较小数转移。暴力版的慢点要落到本题：慢版本会围绕“每个数的一的个数可由去掉最低位一后的较小数转移”使用集合、计数或完整状态表。样例规律是“规律是每个数的答案来自去掉最低位 1 或右移一位后的更小数”，说明哪些信息可以被逐位抵消、压缩或复用；位运算只是编码，不是证明本身。后面的优化只消掉这类重复工作，不能改变答案集合。

\textbf{优化条件。} 本题的关键观察是：dp[x] = dp[x \& (x - 1)] + 1，或 dp[x] = dp[x >> 1] + (x \& 1)。这句话必须能翻译成可维护状态；如果题目条件改变后观察不再成立，就不能继续套原来的写法。

\textbf{相邻状态。} 规律是每个数的答案来自去掉最低位 1 或右移一位后的更小数。把这个特例继续拆开看：每一位或每个掩码位都要对应题目里的一个布尔事实。位运算只是压缩表示；如果两个状态在位置相同但掩码不同，未来能力不同，就不能合并。

\textbf{状态复用。} 把集合状态压成掩码；包含、加入、删除、枚举子集都变成位操作，但每一位的题目含义不能丢。本题落点是：规律是每个数的答案来自去掉最低位 1 或右移一位后的更小数。手算时只改被新输入影响的那一小部分，而不是回头重算完整候选。

\textbf{指针和字段。} 状态字段要能说清：掩码含义、当前位、目标位集合、状态转移和答案读取；负数或高位参与时要说明整数宽度。手算时按这个协议跟踪：拿 n=5 手算，5 的二进制 101，比 4 多一个最低位 1，所以 bits[5]=bits[4]+1；也可用 5\&(4)=4 得到 bits[5]=bits[4]+1。然后按协议复查：手算时把整数写成二进制，并在每一位旁边标注含义。状态转移只允许改变题目动作允许改变的那些位。

\textbf{代码落点。} 移位表达式加括号；使用无符号或足够宽类型；枚举子集时确认空集和全集是否包含。关键代码行必须能逐句翻译成“旧状态怎样变成新状态”，否则就只是模板。

\textbf{正确性。} 证明从逐位独立或子集归纳开始：被压缩到位里的信息足以决定未来转移。

\textbf{边界实验。} 先用n 为零、二的幂、连续区间、表达式优先级做反例检查。实验时覆盖n 为零、二的幂、连续区间、表达式优先级，把数字写成二进制逐位检查。负数、最高位、空掩码和全集掩码都要单独测。

\textbf{复杂度变化。} 暴力用集合或数组保存状态可能需要更多空间和比较成本；位运算通常把每步状态更新压成常数或按位数计算，掩码 DP 仍按状态数增长。

\textbf{迁移追问。} 如果题目改成“它展示位运算和 DP 的结合，而不是逐个数循环数位”，先判断原状态是否仍然覆盖新答案集合，再决定是微调边界、改 value 语义，还是换算法。

\noindent\textbf{LeetCode 371 两整数之和。}

\textbf{题目说明。} LeetCode 371 两整数之和 的题面入口要先还原成具体任务：不用加号时，加法可以拆成无进位和与进位。

\textbf{样例入口。} 拿 5+7 手算，异或得到无进位和 2，按位与左移得到进位 10；继续迭代直到进位为 0。

\textbf{暴力基线。} 慢版本先用集合、数组或布尔表保存状态，逐步执行题目动作。确认每个布尔字段的含义后，才能把它压缩成位，而不是直接写位运算公式。放到本题就是：先按“不用加号时，加法可以拆成无进位和与进位”完整试慢版本；样例里已经能看到“规律是加法拆成无进位和与进位两部分，负数参与时要用足够明确的整数宽度处理”，所以优化前必须先能说清慢版本枚举了哪些候选。

\textbf{暴力过程。} 拿 5+7 手算，异或得到无进位和 2，按位与左移得到进位 10；继续迭代直到进位为 0。

\textbf{重复工作。} 题面模型：不用加号时，加法可以拆成无进位和与进位。暴力版的慢点要落到本题：慢版本会围绕“不用加号时，加法可以拆成无进位和与进位”使用集合、计数或完整状态表。样例规律是“规律是加法拆成无进位和与进位两部分，负数参与时要用足够明确的整数宽度处理”，说明哪些信息可以被逐位抵消、压缩或复用；位运算只是编码，不是证明本身。后面的优化只消掉这类重复工作，不能改变答案集合。

\textbf{优化条件。} 本题的关键观察是：异或得到无进位和，与运算后左移得到进位，迭代直到进位为零。这句话必须能翻译成可维护状态；如果题目条件改变后观察不再成立，就不能继续套原来的写法。

\textbf{相邻状态。} 规律是加法拆成无进位和与进位两部分，负数参与时要用足够明确的整数宽度处理。把这个特例继续拆开看：每一位或每个掩码位都要对应题目里的一个布尔事实。位运算只是压缩表示；如果两个状态在位置相同但掩码不同，未来能力不同，就不能合并。

\textbf{状态复用。} 把集合状态压成掩码；包含、加入、删除、枚举子集都变成位操作，但每一位的题目含义不能丢。本题落点是：规律是加法拆成无进位和与进位两部分，负数参与时要用足够明确的整数宽度处理。手算时只改被新输入影响的那一小部分，而不是回头重算完整候选。

\textbf{指针和字段。} 状态字段要能说清：掩码含义、当前位、目标位集合、状态转移和答案读取；负数或高位参与时要说明整数宽度。手算时按这个协议跟踪：拿 5+7 手算，异或得到无进位和 2，按位与左移得到进位 10；继续迭代直到进位为 0。然后按协议复查：手算时把整数写成二进制，并在每一位旁边标注含义。状态转移只允许改变题目动作允许改变的那些位。

\textbf{代码落点。} 移位表达式加括号；使用无符号或足够宽类型；枚举子集时确认空集和全集是否包含。关键代码行必须能逐句翻译成“旧状态怎样变成新状态”，否则就只是模板。

\textbf{正确性。} 证明从逐位独立或子集归纳开始：被压缩到位里的信息足以决定未来转移。

\textbf{边界实验。} 先用正负数、进位传播、有符号溢出、需要使用无符号中间值做反例检查。实验时覆盖正负数、进位传播、有符号溢出、需要使用无符号中间值，把数字写成二进制逐位检查。负数、最高位、空掩码和全集掩码都要单独测。

\textbf{复杂度变化。} 暴力用集合或数组保存状态可能需要更多空间和比较成本；位运算通常把每步状态更新压成常数或按位数计算，掩码 DP 仍按状态数增长。

\textbf{迁移追问。} 如果题目改成“这是位级算术模拟，和异或抵消类题目的语义不同”，先判断原状态是否仍然覆盖新答案集合，再决定是微调边界、改 value 语义，还是换算法。

\noindent\textbf{LeetCode 372 超级次方。}

\textbf{题目说明。} LeetCode 372 超级次方 的题面入口要先还原成具体任务：大指数以十进制数组给出，不能先转成普通整数。

\textbf{样例入口。} 拿 a=2、b=[1, 0] 手算，目标是 \ensuremath{2^{10}} mod 1337，不能把很长指数转成普通整数。扫描十进制指数时，新答案等于旧答案的十次方乘 a 的当前位次方。

\textbf{暴力基线。} 慢版本先用集合、数组或布尔表保存状态，逐步执行题目动作。确认每个布尔字段的含义后，才能把它压缩成位，而不是直接写位运算公式。放到本题就是：先按“大指数以十进制数组给出，不能先转成普通整数”完整试慢版本；样例里已经能看到“规律是把大指数按十进制位折叠，每步都用快速幂并取模”，所以优化前必须先能说清慢版本枚举了哪些候选。

\textbf{暴力过程。} 拿 a=2、b=[1, 0] 手算，目标是 \ensuremath{2^{10}} mod 1337，不能把很长指数转成普通整数。扫描十进制指数时，新答案等于旧答案的十次方乘 a 的当前位次方。

\textbf{重复工作。} 题面模型：大指数以十进制数组给出，不能先转成普通整数。暴力版的慢点要落到本题：慢版本会围绕“大指数以十进制数组给出，不能先转成普通整数”使用集合、计数或完整状态表。样例规律是“规律是把大指数按十进制位折叠，每步都用快速幂并取模”，说明哪些信息可以被逐位抵消、压缩或复用；位运算只是编码，不是证明本身。后面的优化只消掉这类重复工作，不能改变答案集合。

\textbf{优化条件。} 本题的关键观察是：利用 \ensuremath{(a^10)^{rest}} 和最后一位指数，把指数数组逐位折叠，并在每步取模。这句话必须能翻译成可维护状态；如果题目条件改变后观察不再成立，就不能继续套原来的写法。

\textbf{相邻状态。} 规律是把大指数按十进制位折叠，每步都用快速幂并取模。把这个特例继续拆开看：每一位或每个掩码位都要对应题目里的一个布尔事实。位运算只是压缩表示；如果两个状态在位置相同但掩码不同，未来能力不同，就不能合并。

\textbf{状态复用。} 把集合状态压成掩码；包含、加入、删除、枚举子集都变成位操作，但每一位的题目含义不能丢。本题落点是：规律是把大指数按十进制位折叠，每步都用快速幂并取模。手算时只改被新输入影响的那一小部分，而不是回头重算完整候选。

\textbf{指针和字段。} 状态字段要能说清：掩码含义、当前位、目标位集合、状态转移和答案读取；负数或高位参与时要说明整数宽度。手算时按这个协议跟踪：拿 a=2、b=[1, 0] 手算，目标是 \ensuremath{2^{10}} mod 1337，不能把很长指数转成普通整数。扫描十进制指数时，新答案等于旧答案的十次方乘 a 的当前位次方。然后按协议复查：手算时把整数写成二进制，并在每一位旁边标注含义。状态转移只允许改变题目动作允许改变的那些位。

\textbf{代码落点。} 移位表达式加括号；使用无符号或足够宽类型；枚举子集时确认空集和全集是否包含。关键代码行必须能逐句翻译成“旧状态怎样变成新状态”，否则就只是模板。

\textbf{正确性。} 证明从逐位独立或子集归纳开始：被压缩到位里的信息足以决定未来转移。

\textbf{边界实验。} 先用指数为空、a 大于模数、指数含零、递归深度、模乘溢出做反例检查。实验时覆盖指数为空、a 大于模数、指数含零、递归深度、模乘溢出，把数字写成二进制逐位检查。负数、最高位、空掩码和全集掩码都要单独测。

\textbf{复杂度变化。} 暴力用集合或数组保存状态可能需要更多空间和比较成本；位运算通常把每步状态更新压成常数或按位数计算，掩码 DP 仍按状态数增长。

\textbf{迁移追问。} 如果题目改成“若模数变化或要求精确值，取模快速幂的状态语义要重写”，先判断原状态是否仍然覆盖新答案集合，再决定是微调边界、改 value 语义，还是换算法。

\noindent\textbf{LeetCode 393 UTF-8 编码验证。}

\textbf{题目说明。} LeetCode 393 UTF-8 编码验证 的题面入口要先还原成具体任务：每个字节的高位模式决定一个字符需要的后续字节数量。

\textbf{样例入口。} 拿 [197, 130, 1] 手算，197 的高位模式表示还需要 1 个 continuation 字节，130 满足 10xxxxxx，之后 1 是单字节。

\textbf{暴力基线。} 慢版本先用集合、数组或布尔表保存状态，逐步执行题目动作。确认每个布尔字段的含义后，才能把它压缩成位，而不是直接写位运算公式。放到本题就是：先按“每个字节的高位模式决定一个字符需要的后续字节数量”完整试慢版本；样例里已经能看到“规律是状态保存剩余后续字节数量，非法首字节或后续字节不足都失败”，所以优化前必须先能说清慢版本枚举了哪些候选。

\textbf{暴力过程。} 拿 [197, 130, 1] 手算，197 的高位模式表示还需要 1 个 continuation 字节，130 满足 10xxxxxx，之后 1 是单字节。

\textbf{重复工作。} 题面模型：每个字节的高位模式决定一个字符需要的后续字节数量。暴力版的慢点要落到本题：慢版本会围绕“每个字节的高位模式决定一个字符需要的后续字节数量”使用集合、计数或完整状态表。样例规律是“规律是状态保存剩余后续字节数量，非法首字节或后续字节不足都失败”，说明哪些信息可以被逐位抵消、压缩或复用；位运算只是编码，不是证明本身。后面的优化只消掉这类重复工作，不能改变答案集合。

\textbf{优化条件。} 本题的关键观察是：用位掩码判断首字节类型，并维护还需要多少个 continuation 字节。这句话必须能翻译成可维护状态；如果题目条件改变后观察不再成立，就不能继续套原来的写法。

\textbf{相邻状态。} 规律是状态保存剩余后续字节数量，非法首字节或后续字节不足都失败。把这个特例继续拆开看：每一位或每个掩码位都要对应题目里的一个布尔事实。位运算只是压缩表示；如果两个状态在位置相同但掩码不同，未来能力不同，就不能合并。

\textbf{状态复用。} 把集合状态压成掩码；包含、加入、删除、枚举子集都变成位操作，但每一位的题目含义不能丢。本题落点是：规律是状态保存剩余后续字节数量，非法首字节或后续字节不足都失败。手算时只改被新输入影响的那一小部分，而不是回头重算完整候选。

\textbf{指针和字段。} 状态字段要能说清：掩码含义、当前位、目标位集合、状态转移和答案读取；负数或高位参与时要说明整数宽度。手算时按这个协议跟踪：拿 [197, 130, 1] 手算，197 的高位模式表示还需要 1 个 continuation 字节，130 满足 10xxxxxx，之后 1 是单字节。然后按协议复查：手算时把整数写成二进制，并在每一位旁边标注含义。状态转移只允许改变题目动作允许改变的那些位。

\textbf{代码落点。} 移位表达式加括号；使用无符号或足够宽类型；枚举子集时确认空集和全集是否包含。关键代码行必须能逐句翻译成“旧状态怎样变成新状态”，否则就只是模板。

\textbf{正确性。} 证明从逐位独立或子集归纳开始：被压缩到位里的信息足以决定未来转移。

\textbf{边界实验。} 先用非法首字节、后续字节不足、后续字节格式错误、单字节字符做反例检查。实验时覆盖非法首字节、后续字节不足、后续字节格式错误、单字节字符，把数字写成二进制逐位检查。负数、最高位、空掩码和全集掩码都要单独测。

\textbf{复杂度变化。} 暴力用集合或数组保存状态可能需要更多空间和比较成本；位运算通常把每步状态更新压成常数或按位数计算，掩码 DP 仍按状态数增长。

\textbf{迁移追问。} 如果题目改成“这是位模式解析题，重点是协议规则和状态机”，先判断原状态是否仍然覆盖新答案集合，再决定是微调边界、改 value 语义，还是换算法。

\noindent\textbf{LeetCode 421 数组中两个数的最大异或值。}

\textbf{题目说明。} LeetCode 421 数组中两个数的最大异或值 的题面入口要先还原成具体任务：最大异或值由高位到低位尽量取一决定。

\textbf{样例入口。} 拿 [3, 10, 5, 25, 2, 8] 手算，25 和 5 的异或为 28。逐位构造答案时，先假设当前位能为 1，再用哈希集合检查是否存在两个前缀配对。

\textbf{暴力基线。} 慢版本先用集合、数组或布尔表保存状态，逐步执行题目动作。确认每个布尔字段的含义后，才能把它压缩成位，而不是直接写位运算公式。放到本题就是：先按“最大异或值由高位到低位尽量取一决定”完整试慢版本；样例里已经能看到“规律是从高位到低位贪心验证前缀，不需要枚举所有数对”，所以优化前必须先能说清慢版本枚举了哪些候选。

\textbf{暴力过程。} 拿 [3, 10, 5, 25, 2, 8] 手算，25 和 5 的异或为 28。逐位构造答案时，先假设当前位能为 1，再用哈希集合检查是否存在两个前缀配对。

\textbf{重复工作。} 题面模型：最大异或值由高位到低位尽量取一决定。暴力版的慢点要落到本题：慢版本会围绕“最大异或值由高位到低位尽量取一决定”使用集合、计数或完整状态表。样例规律是“规律是从高位到低位贪心验证前缀，不需要枚举所有数对”，说明哪些信息可以被逐位抵消、压缩或复用；位运算只是编码，不是证明本身。后面的优化只消掉这类重复工作，不能改变答案集合。

\textbf{优化条件。} 本题的关键观察是：逐位尝试把当前前缀答案设为一，用哈希集合验证是否存在两个前缀配对。这句话必须能翻译成可维护状态；如果题目条件改变后观察不再成立，就不能继续套原来的写法。

\textbf{相邻状态。} 规律是从高位到低位贪心验证前缀，不需要枚举所有数对。把这个特例继续拆开看：每一位或每个掩码位都要对应题目里的一个布尔事实。位运算只是压缩表示；如果两个状态在位置相同但掩码不同，未来能力不同，就不能合并。

\textbf{状态复用。} 把集合状态压成掩码；包含、加入、删除、枚举子集都变成位操作，但每一位的题目含义不能丢。本题落点是：规律是从高位到低位贪心验证前缀，不需要枚举所有数对。手算时只改被新输入影响的那一小部分，而不是回头重算完整候选。

\textbf{指针和字段。} 状态字段要能说清：掩码含义、当前位、目标位集合、状态转移和答案读取；负数或高位参与时要说明整数宽度。手算时按这个协议跟踪：拿 [3, 10, 5, 25, 2, 8] 手算，25 和 5 的异或为 28。逐位构造答案时，先假设当前位能为 1，再用哈希集合检查是否存在两个前缀配对。然后按协议复查：手算时把整数写成二进制，并在每一位旁边标注含义。状态转移只允许改变题目动作允许改变的那些位。

\textbf{代码落点。} 移位表达式加括号；使用无符号或足够宽类型；枚举子集时确认空集和全集是否包含。关键代码行必须能逐句翻译成“旧状态怎样变成新状态”，否则就只是模板。

\textbf{正确性。} 证明从逐位独立或子集归纳开始：被压缩到位里的信息足以决定未来转移。

\textbf{边界实验。} 先用零、重复数、最高位、负数语义做反例检查。实验时覆盖零、重复数、最高位、负数语义，把数字写成二进制逐位检查。负数、最高位、空掩码和全集掩码都要单独测。

\textbf{复杂度变化。} 暴力用集合或数组保存状态可能需要更多空间和比较成本；位运算通常把每步状态更新压成常数或按位数计算，掩码 DP 仍按状态数增长。

\textbf{迁移追问。} 如果题目改成“也可以建二进制 Trie，每个数尽量走相反位”，先判断原状态是否仍然覆盖新答案集合，再决定是微调边界、改 value 语义，还是换算法。

\noindent\textbf{LeetCode 1755 最接近目标值的子序列和。}

\textbf{题目说明。} LeetCode 1755 最接近目标值的子序列和 的题面入口要先还原成具体任务：完整子序列选择是 \ensuremath{2^{n}}，拆成左右两半后各自枚举子集和。

\textbf{样例入口。} 拿 nums=[5, -7, 3, 5]、goal=6 手算，完整枚举 16 个子序列可找到 -7+3+5+5=6。折半后左半 [5, -7] 生成 0、5、-7、-2，右半 [3, 5] 生成 0、3、5、8；对每个左半和二分找最接近 goal-left 的右半和。

\textbf{暴力基线。} 慢版本先用集合、数组或布尔表保存状态，逐步执行题目动作。确认每个布尔字段的含义后，才能把它压缩成位，而不是直接写位运算公式。放到本题就是：先按“完整子序列选择是 \ensuremath{2^{n}}，拆成左右两半后各自枚举子集和”完整试慢版本；样例里已经能看到“规律是把 \ensuremath{2^{n}} 拆成两个 \ensuremath{2^{n/2}} 的子集和集合再配对”，所以优化前必须先能说清慢版本枚举了哪些候选。

\textbf{暴力过程。} 拿 nums=[5, -7, 3, 5]、goal=6 手算，完整枚举 16 个子序列可找到 -7+3+5+5=6。折半后左半 [5, -7] 生成 0、5、-7、-2，右半 [3, 5] 生成 0、3、5、8；对每个左半和二分找最接近 goal-left 的右半和。

\textbf{重复工作。} 题面模型：完整子序列选择是 \ensuremath{2^{n}}，拆成左右两半后各自枚举子集和。暴力版的慢点要落到本题：慢版本会围绕“完整子序列选择是 \ensuremath{2^{n}}，拆成左右两半后各自枚举子集和”使用集合、计数或完整状态表。样例规律是“规律是把 \ensuremath{2^{n}} 拆成两个 \ensuremath{2^{n/2}} 的子集和集合再配对”，说明哪些信息可以被逐位抵消、压缩或复用；位运算只是编码，不是证明本身。后面的优化只消掉这类重复工作，不能改变答案集合。

\textbf{优化条件。} 本题的关键观察是：分别生成两半所有子集和，排序一侧，对另一侧每个和用二分找最接近 goal 的补数。这句话必须能翻译成可维护状态；如果题目条件改变后观察不再成立，就不能继续套原来的写法。

\textbf{相邻状态。} 规律是把 \ensuremath{2^{n}} 拆成两个 \ensuremath{2^{n/2}} 的子集和集合再配对。把这个特例继续拆开看：每一位或每个掩码位都要对应题目里的一个布尔事实。位运算只是压缩表示；如果两个状态在位置相同但掩码不同，未来能力不同，就不能合并。

\textbf{状态复用。} 把集合状态压成掩码；包含、加入、删除、枚举子集都变成位操作，但每一位的题目含义不能丢。本题落点是：规律是把 \ensuremath{2^{n}} 拆成两个 \ensuremath{2^{n/2}} 的子集和集合再配对。手算时只改被新输入影响的那一小部分，而不是回头重算完整候选。

\textbf{指针和字段。} 状态字段要能说清：掩码含义、当前位、目标位集合、状态转移和答案读取；负数或高位参与时要说明整数宽度。手算时按这个协议跟踪：拿 nums=[5, -7, 3, 5]、goal=6 手算，完整枚举 16 个子序列可找到 -7+3+5+5=6。折半后左半 [5, -7] 生成 0、5、-7、-2，右半 [3, 5] 生成 0、3、5、8；对每个左半和二分找最接近 goal-left 的右半和。然后按协议复查：手算时把整数写成二进制，并在每一位旁边标注含义。状态转移只允许改变题目动作允许改变的那些位。

\textbf{代码落点。} 移位表达式加括号；使用无符号或足够宽类型；枚举子集时确认空集和全集是否包含。关键代码行必须能逐句翻译成“旧状态怎样变成新状态”，否则就只是模板。

\textbf{正确性。} 证明从逐位独立或子集归纳开始：被压缩到位里的信息足以决定未来转移。

\textbf{边界实验。} 先用负数、空子序列、目标为负、答案为零提前返回、两半大小做反例检查。实验时覆盖负数、空子序列、目标为负、答案为零提前返回、两半大小，把数字写成二进制逐位检查。负数、最高位、空掩码和全集掩码都要单独测。

\textbf{复杂度变化。} 暴力用集合或数组保存状态可能需要更多空间和比较成本；位运算通常把每步状态更新压成常数或按位数计算，掩码 DP 仍按状态数增长。

\textbf{迁移追问。} 如果题目改成“若要求方案数而非最接近差值，要改成计数配对”，先判断原状态是否仍然覆盖新答案集合，再决定是微调边界、改 value 语义，还是换算法。

\subsection{实验复盘}

追问通常会问每一位代表什么、负数和移位是否安全、状态相同但资源不同能否合并。请准备说明位运算只是编码，真正关键仍是状态语义。

复盘本节时先从 LeetCode 29 两数相除、LeetCode 50 Pow(x, n)、LeetCode 136 只出现一次的数字 开始：每道题都先手写一个能覆盖全部候选的暴力版，再指出优化结构具体保存了哪一段历史信息，最后用相邻案例比较为什么不能互换解法。
